\section{Overview}\label{sec:map}
The paper asks when equal server metadata at two captures proves that a
resource stayed unchanged in between, and when captures of several resources
prove that their observed contents were current together. This supplement
gives the proofs in the paper's order (Sections~\ref{sec:continuity}
to~\ref{sec:instants}), then the evidence: platform semantics, local
executions and the documentation audit (Section~\ref{sec:platforms}), the
edit-history replay (Section~\ref{sec:replay}), the controlled HTTP executions
(Section~\ref{sec:http}) and reproduction (Section~\ref{sec:reproduction}).
Table~\ref{tab:map} lists where each claim is proved or measured.

Four kinds of evidence are kept apart throughout. A proof holds for the
stated model. A finite enumeration checks an implementation against an
independent brute-force oracle; it does not replace a proof. A local execution
records what one pinned implementation did in one run on loopback; it is not
a cloud measurement and estimates no prevalence. Vendor documentation states
intended behavior, and the paper's claims about cloud services rest on it
alone: no AWS, Google Cloud or Azure service was measured.

\begin{table*}[t]
\caption{Where each claim of the paper is proved or measured.}
\label{tab:map}
\small
\begin{tabularx}{\textwidth}{@{}lXlX@{}}
\toprule
Paper & Claim & Here & Evidence, code and data \\
\midrule
Theorem~1 & Continuity criterion and its certificates & \S\ref{sec:continuity} &
Proof; product search and invariant checker in \code{src/history_contracts.py};
13 constructed models and 37{,}888 exhaustive contracts against an
independent oracle \\
Proposition~2 & Sound fields, no finite path-preserving repair, identity
non-reactivation &
\S\ref{sec:fields} & Proof; \code{repair()} and \code{refine()}; \sRepairChecks{}
contracts against all labelings, \sRefinementChecks{} finite refinements,
\sEpochChecks{} bounded S3 epochs \\
\S2, \S4 & Documented operations of Amazon S3 and other platforms; local probes; audit &
\S\ref{sec:platforms} & Vendor documentation; \nProbeRequests{} loopback
requests on nine implementations; two MinIO executions (51 and
\nProfileRequests{} requests); \nAuditIdentifiers{} identifiers from
\nAuditSources{} sources \\
Theorem~3, Corollary~4 & Exact test, schedules, fixed instants &
\S\ref{sec:instants} & Proof; \code{src/certification.py}; schedules and window
sets against brute force \\
Proposition~5 & Worst-case sweep, coNP-completeness & \S\ref{sec:instants} &
Proof; sweep against the reference semantics; \sCnfChecks{} CNF instances
of the reduction \\
Figure~1 & Recorded MinIO execution; revisit order & \S\ref{sec:platforms},
\S\ref{sec:instants} & \code{figure-artifacts/05-certification-rule} \\
Table~1 & Identifiers across a change and its undo & \S\ref{sec:platforms} &
\code{table-artifacts/01-identity-probes} \\
Figure~2, \S4 & Edit-history replay & \S\ref{sec:replay} & Protocol committed
before collection; two labelled post hoc addenda; \code{experiments/revision_bundles},
\code{figure-artifacts/06-revision-histories} \\
\S1, \S2 & Validators from restored modification times repeat (HTTP controls) &
\S\ref{sec:http} & Six configurations on macOS and on Linux;
\code{scripts/run_stack_experiments.py} \\
\bottomrule
\end{tabularx}
\end{table*}

The supplement covers only what the paper uses. Results of earlier rounds that
the paper no longer prints (finite-capture ambiguity, the unbounded obstruction
order, label-alphabet synthesis, a six-page source-history study) are kept,
unchanged, in \code{research/supplement-earlier-rounds.md}; nothing here
depends on them.
